\section*{Chapter \ref{ch:b2:orbital-forcing} --- Orbital Forcing and the Ice Ages}

\begin{solution}{exo:b2:orbital-forcing:1}
(a) $1.0167/0.9833 = 1.0340$: the aphelion is \qty{3.4}{\percent} farther than the perihelion.
(b) $a(1 + e) - a(1 - e) = 2ea = 0.0334\,a$, \qty{3.34}{\percent} of $a$, close to the 3.4 per cent
usually quoted. (c) Sunlight falls off as the inverse square of the distance: $1.0340^2 = 1.069$,
about \qty{7}{\percent} more sunlight at perihelion.
\end{solution}

\begin{solution}{exo:b2:orbital-forcing:2}
(a) Today $0.017236 \times \sin\ang{101.37} = 0.017236 \times 0.9804 = 0.01690$; \num{10000} years ago
$0.019971 \times \sin\ang{293.86} = 0.019971 \times (-0.9145) = -0.01826$; \num{6000} years ago
$0.019249 \times \sin\ang{359.99} \approx 0$, the values of the tables. (b) The Sun's longitude at
perihelion is $\varpi + \ang{180}$: \ang{281.37} today, \ang{473.86} $-$ \ang{360} $=$ \ang{113.86}
\num{10000} years ago, \ang{179.99} \num{6000} years ago. (c) \num{10000} years ago the perihelion
fell \ang{24} after the June solstice, the nearest of the three to it, in the northern summer;
\num{6000} years ago it fell at the September equinox, where the \omterm{def:b2:orbital-forcing:elements}{climatic precession} is zero.
\end{solution}

\begin{solution}{exo:b2:orbital-forcing:3}
(a) The frequencies add: $1/25.77 + 1/112 = 0.03880 + 0.00893 = \qty{0.04773}{kyr^{-1}}$, a period
of $1/0.04773 = 20.9$ thousand years. (b) The result falls within the range of the main periods,
19.0 to 23.7 thousand years. A single number cannot reproduce them: the turning of the orbit
is itself the sum of several motions with different frequencies, each of which combines with the
precession of the axis to give one period of the \omterm{def:b2:orbital-forcing:elements}{climatic precession}; the average is the
\num{23000} years usually quoted.
\end{solution}

\begin{solution}{exo:b2:orbital-forcing:4}
$\sqrt{1 - 0.0167^2} = 0.99986$. For \ang{22.1}: $1360 \times 0.37622/(\pi \times 0.99986) =
\qty{162.9}{W.m^{-2}}$; for \ang{24.5}: $1360 \times 0.41469/(\pi \times 0.99986) =
\qty{179.5}{W.m^{-2}}$. The ratio is $\sin\ang{24.5}/\sin\ang{22.1} = 1.102$: the annual mean at the
pole grows by \qty{10.2}{\percent}, about \qty{16.7}{W.m^{-2}}, between the two extremes.
\end{solution}

\begin{solution}{exo:b2:orbital-forcing:5}
(a) At the June solstice $(a/r)^2 = \bigl((1 - e\sin\varpi)/(1 - e^2)\bigr)^2$. Today $e\sin\varpi
= 0.016898$ and $(0.983102/0.999703)^2 = 0.9671$, against $1 - 0.033796 = 0.9662$;
\num{116000} years ago $e\sin\varpi = 0.043939$ and $(0.956061/0.998065)^2 = 0.9176$, against $1 -
0.087878 = 0.9121$. (b) At the December solstice $(a/r)^2 = \bigl((1 + e\sin\varpi)/(1 -
e^2)\bigr)^2$: today $1.0347$ against $1.0338$, \num{116000} years ago $1.0940$ against $1.0879$.
(c) The neglected terms are of order $e^2$: $(1 - x)^2/(1 - e^2)^2 \approx 1 - 2x + x^2 + 2e^2$
with $x = e\sin\varpi$, here about $3e^2$. They are $3 \times 0.0172^2 = 0.0009$ today and $3
\times 0.0440^2 = 0.0058$ \num{116000} years ago, close to the differences found, $0.0009$ and
$0.0055$: the approximation is six times worse when the eccentricity is two and a half times
larger.
\end{solution}

\begin{solution}{exo:b2:orbital-forcing:6}
At the June solstice $\delta = \varepsilon = \ang{23.446}$. $\cos h_0 = -\tan\ang{65}\tan\ang{23.446}
= -2.14451 \times 0.43369 = -0.93006$, so $h_0 = \ang{158.44} = \qty{2.7654}{rad}$. The day lasts
$2 \times 158.44/15 = 21.13$~hours (the Earth turns \ang{15} an hour). With $\sin\ang{65} = 0.90631$,
$\cos\ang{65} = 0.42262$, $\sin\delta = 0.39788$, $\cos\delta = 0.91744$ and $\sin h_0 = 0.36742$,
the bracket is $2.7654 \times 0.90631 \times 0.39788 + 0.42262 \times 0.91744 \times 0.36742 = 0.99720
+ 0.14246 = 1.13966$. At the mean distance $\bar Q = (1360/\pi) \times 1.13966 = 432.90 \times 1.13966
= \qty{493.4}{W.m^{-2}}$; with the distance factor of \cref{exo:b2:orbital-forcing:5}, $493.4 \times
0.9671 = \qty{477.1}{W.m^{-2}}$. As a check, the same steps at \ang{60} N give \qty{475.96}{W.m^{-2}}, where the tables give \num{475.95}.
\end{solution}

\begin{solution}{exo:b2:orbital-forcing:7}
(a) The March equinox is $\lambda = 0$, the September equinox $\lambda = \pi$. With $\sqrt{1 - e^2}
\approx 1$, the duration is $\frac{T}{2\pi}\int_0^{\pi}\bigl(1 + 2e\cos(\lambda -
\varpi)\bigr)\,\mathrm d\lambda = \frac{T}{2\pi}\bigl(\pi + 2e[\sin(\lambda - \varpi)]_0^{\pi}\bigr)
= \frac{T}{2\pi}(\pi + 4e\sin\varpi) = \frac T2 + \frac{2T}{\pi}e\sin\varpi$. (b) The other half-year
is shorter by the same amount, so the difference is $(4T/\pi)\,e\sin\varpi = (4 \times
365.2422/\pi) \times 0.016898 = 465.04 \times 0.016898 = 7.86$~days. (c) Spring and summer together
exceed autumn and winter by $4.5 + 3 = 7.5$~days; the two differences quoted are rounded, and the
agreement is within that rounding.
\end{solution}

\begin{solution}{exo:b2:orbital-forcing:8}
(a) For $e = 0.0046$: $\sqrt{1 - e^2} = 0.999989$ and $340/0.999989 = \qty{340.004}{W.m^{-2}}$. For
$e = 0.0565$: $\sqrt{1 - 0.003192} = 0.998403$ and $340/0.998403 = \qty{340.544}{W.m^{-2}}$. (b) The
difference is \qty{0.54}{W.m^{-2}}, \qty{0.16}{\percent}. (c) The mid-July insolation at \ang{65} N
swung by $497 - 388 = \qty{109}{W.m^{-2}}$, about \qty{25}{\percent} of its mean and two hundred
times the change of the global annual mean: the orbit acts on climate by moving sunlight between
seasons and latitudes, not by changing the total.
\end{solution}

\begin{solution}{exo:b2:orbital-forcing:9}
With $\omega = 2\pi/P$ and $\tau = \qty{10}{kyr}$: for $P = 23$, $\omega\tau = 2.732$, a relative
amplitude $1/\sqrt{1 + 7.464} = 0.344$ and a lag $\arctan(2.732)/\omega = 1.2199/0.2732 = 4.47$~kyr;
for $P = 41$, $\omega\tau = 1.532$, amplitude $0.546$, lag $0.9926/0.1532 = 6.48$~kyr; for $P =
100$, $\omega\tau = 0.628$, amplitude $0.847$, lag $0.5610/0.0628 = 8.93$~kyr. The ice sheet
favours the slow forcing: it follows \qty{85}{\percent} of a \num{100000}-year cycle but only a
third of a precession cycle, since it has too little time to grow or shrink before a fast forcing
reverses. It acts as a low-pass filter.
\end{solution}

\begin{solution}{exo:b2:orbital-forcing:10}
(a) The equation is linear: if $V_1$ answers $A_1\sin\omega_1t$ and $V_2$ answers
$A_2\sin\omega_2t$, $V_1 + V_2$ answers their sum. By \cref{prop:b2:orbital-forcing:lag} $V_1$ and
$V_2$ are sinusoids of frequencies $\omega_1$ and $\omega_2$ (once the transient has died away), so
the response holds no other frequency. (b) With $V \approx B_1\sin(\omega_1t - \phi_1) +
B_2\sin(\omega_2t - \phi_2)$,
\[
\begin{aligned}
  V^2 &= \tfrac12(B_1^2 + B_2^2) - \tfrac12B_1^2\cos(2\omega_1t - 2\phi_1) -
  \tfrac12B_2^2\cos(2\omega_2t - 2\phi_2)\\
  &\quad + B_1B_2\bigl[\cos\bigl((\omega_1 - \omega_2)t - \phi_1 + \phi_2\bigr) -
  \cos\bigl((\omega_1 + \omega_2)t - \phi_1 - \phi_2\bigr)\bigr],
\end{aligned}
\]
so the term $kV^2$ forces the equation at the difference and sum frequencies, and at the double
frequencies. (c) $1/(1/19.0 - 1/23.7) = 1/(0.05263 - 0.04219) = 95.8$ thousand years; the pair 19.0
and 22.4 gives 125 thousand years. The difference frequencies of the precession terms fall in the
eccentricity band, as they must, since the \omterm{def:b2:orbital-forcing:elements}{climatic precession} is $e\sin\varpi$. A non-linear
ice sheet could therefore draw a \num{100000}-year rhythm out of the precession; a linear one
cannot. The \num{100000}-year cycles are evidence of non-linear behaviour, though this calculation
alone does not show which.
\end{solution}

\begin{solution}{exo:b2:orbital-forcing:11}
(a) As in the proof of \cref{prop:b2:orbital-forcing:annual}, $(a/r)^2\,\mathrm dt =
T\,\mathrm d\lambda/(2\pi\sqrt{1 - e^2})$, and $\sin\delta = \sin\varepsilon\sin\lambda$ is positive
for $0 < \lambda < \pi$. So $E = \frac{S_0T\sin\varepsilon}{2\pi\sqrt{1 - e^2}}\int_0^{\pi}\sin\lambda
\,\mathrm d\lambda = S_0T\sin\varepsilon/(\pi\sqrt{1 - e^2})$. (b) $T = 365.25 \times 86400 =
\qty{3.1558e7}{s}$, $S_0T/\pi = \qty{1.3661e10}{J.m^{-2}}$; for \ang{22.1}, $E = 1.3661 \times 10^{10}
\times 0.37622/0.99986 = \qty{5.14}{GJ.m^{-2}}$; for \ang{24.5}, $\qty{5.67}{GJ.m^{-2}}$,
\qty{10}{\percent} more. (c) The precession fixes how the same energy is spread in time: a summer
at perihelion is shorter and more intense, a summer at aphelion longer and weaker. Melting starts
only above \qty{0}{\celsius}, so a short intense summer can melt more than a long cool one
receiving the same total.
\end{solution}

\begin{solution}{exo:b2:orbital-forcing:12}
(a) $3.93\log_2(279.4/189.7) = 3.93\log_2 1.4729 = 3.93 \times 0.5586 = \qty{2.20}{W.m^{-2}}$. (b) The
smallest rise is from 194 to \qty{266}{ppm}: $3.93\log_2 1.3711 = 3.93 \times 0.4554 =
\qty{1.79}{W.m^{-2}}$; the largest from 188 to \qty{282}{ppm}: $3.93\log_2 1.5 = 3.93 \times 0.5850 =
\qty{2.30}{W.m^{-2}}$. (c) The forcing of carbon dioxide is four times the largest change of the
global annual insolation, \qty{0.54}{W.m^{-2}}, and it acts everywhere and all year; the summer
insolation at \ang{65} N swings by about \qty{109}{W.m^{-2}}, but at one latitude and one season. The
orbit is the trigger, through the summers of the northern ice sheets; carbon dioxide, with the
\omterm{def:b2:precambrian-earth:ice-albedo}{ice--albedo feedback}, is an amplifier that spreads the change over the globe.
\end{solution}

\begin{solution}{pb:b2:orbital-forcing:1}
\textbf{1.} Then $0.0440 \times \sin\ang{92.7} = 0.0440 \times 0.99889 = 0.04395$; now $0.0172 \times
\sin\ang{101.4} = 0.0172 \times 0.98027 = 0.01686$, nearly three times less.
\textbf{2.} Then $1 - 0.0440 = 0.956$ and $1 + 0.0440 = 1.044$; now $0.9828$ and $1.0172$.
\textbf{3.} At $\lambda = \varpi + \ang{180}$: then \ang{272.7}, now \ang{281.4}, both near the
December solstice ($\lambda = \ang{270}$); then the Earth passed perihelion less than three days
after it, and aphelion at the June solstice.
\textbf{4.} Then $(1.044/0.956)^2 = 1.0921^2 = 1.193$, now $(1.0172/0.9828)^2 = 1.0350^2 = 1.071$:
\qty{19}{\percent} more sunlight at perihelion then, \qty{7}{\percent} now.
\textbf{5.} Smaller, \ang{22.52} against \ang{23.45}: the summer sun stood lower at high latitudes,
and high latitudes received less sunlight in summer and over the year.
\textbf{6.} Then $340/\sqrt{1 - 0.001936} = 340/0.999032 = \qty{340.33}{W.m^{-2}}$; now $340/0.999852
= \qty{340.05}{W.m^{-2}}$. The difference, \qty{0.28}{W.m^{-2}} or less than one part in a thousand,
is insignificant.
\textbf{7.} At $\lambda = \ang{90}$, $\cos(\ang{90} - \varpi) = \sin\varpi$ in
\cref{prop:b2:orbital-forcing:distance}.
\textbf{8.} Then $(1 - 0.001936)/(1 - 0.04395) = 0.998064/0.956049 = 1.04395$; now $(1 -
0.000296)/(1 - 0.01686) = 0.999704/0.983139 = 1.01685$.
\textbf{9.} Then $1/1.04395^2 = 0.91758$; now $1/1.01685^2 = 0.96714$.
\textbf{10.} Then $1 - 2 \times 0.04395 = 0.91210$; now $1 - 2 \times 0.01686 = 0.96628$: the
approximation errs by $0.0055$ then, $0.0009$ now.
\textbf{11.} Then $1 - 0.91758 = \qty{8.2}{\percent}$ weaker, now \qty{3.3}{\percent}.
\textbf{12.} At the solstice $\delta = \varepsilon$: \ang{22.52} then, \ang{23.45} now.
\textbf{13.} Then $\cos h_0 = -2.14451 \times 0.41462 = -0.88916$, $h_0 = \ang{152.77} =
\qty{2.6663}{rad}$; now $\cos h_0 = -2.14451 \times 0.43378 = -0.93023$, $h_0 = \ang{158.47} =
\qty{2.7659}{rad}$.
\textbf{14.} $2h_0/\ang{15}$ per hour: then $20.37$~hours, now $21.13$~hours.
\textbf{15.} Then $2.6663 \times 0.90631 \times 0.38301 + 0.42262 \times 0.92375 \times 0.45759 =
0.92554 + 0.17864 = 1.10418$ (\num{1.10417} unrounded), and $432.90 \times 1.10417 =
\qty{478.0}{W.m^{-2}}$; now $2.7659 \times 0.90631 \times 0.39795 + 0.42262 \times 0.91741 \times
0.36697 = 0.99757 + 0.14228 = 1.13985$ (\num{1.13982} unrounded), and $432.90 \times 1.13982 = \qty{493.4}{W.m^{-2}}$.
\textbf{16.} Then $478.0 \times 0.91758 = \qty{438.6}{W.m^{-2}}$; now $493.4 \times 0.96714 =
\qty{477.2}{W.m^{-2}}$, \qty{38.6}{W.m^{-2}} more than then.
\textbf{17.} In mid-July the Sun's longitude is \ang{120}: $\sin\delta = \sin\varepsilon\sin\ang{120}$,
a declination of about \ang{20} instead of \ang{23}, so a shorter day and a lower sun: at
\ang{65} N the daily-mean insolation is greatest near the solstice, when the declination is largest.
\textbf{18.} With today's distance, the change of obliquity alone costs $(493.4 - 478.0) \times
0.96714 = \qty{14.9}{W.m^{-2}}$; with the old obliquity, the distance alone costs $478.0 \times
(0.96714 - 0.91758) = \qty{23.7}{W.m^{-2}}$. The two add to \qty{38.6}{W.m^{-2}}; the distance
counted for about three fifths.
\textbf{19.} After a sudden change the ice sheet has achieved $1 - \mathrm e^{-t/\tau}$ of it:
$1 - \mathrm e^{-0.5} = 0.39$ after \qty{5}{kyr}, $1 - \mathrm e^{-1} = 0.63$ after \qty{10}{kyr},
$1 - \mathrm e^{-2} = 0.86$ after \qty{20}{kyr}.
\textbf{20.} $\omega\tau = 2\pi \times 10/23 = 2.732$: the ice sheet follows $1/\sqrt{1 + 2.732^2} =
0.34$ of the forcing, $\arctan(2.732)/\omega = 4.5$~kyr late.
\textbf{21.} Half a precession cycle, some 11 to \qty{12}{kyr}, is about one response time: before
the ice sheet has covered two thirds of the way, the insolation turns round and pulls it back.
\textbf{22.} In the linear model the ice grows and shrinks at the same rate for forcings of the same
size: a rising and a falling forcing give mirror-image responses. A slow build-up followed by a fast
collapse needs a response that depends on the state or the direction of change, that is a
non-linear one.
\textbf{23.} Carbon dioxide rose by $3.93\log_2(266/194) = 1.8$ to $3.93\log_2(282/188) =
\qty{2.3}{W.m^{-2}}$ from glacial to \omterm{def:b2:orbital-forcing:glacials}{interglacial}: a forcing over the whole globe and the whole
year. Released when the Atlantic circulation weakens, it carries the change of the northern
summers to the whole planet, and, with the albedo of the shrinking ice, amplifies it.
\textbf{24.} When the last \omterm{def:b2:orbital-forcing:glacials}{interglacial} ended, the June-solstice daily-mean insolation at
\ang{65} N was about \textbf{\qty{439}{W.m^{-2}}}, \qty{39}{W.m^{-2}} less than today and its lowest
value of the cycle.
\end{solution}
